\section{Related Work}
\label{sec:related}

\paragraph{From Explicit to Latent Chain-of-Thought.}
Chain-of-Thought has significantly advanced the reasoning capabilities by enabling the decomposition of complex problems into explicit intermediate steps~\citep{wei2022chain}. Early research~\citep{goyal2023think, wang2023guiding, chen-etal-2025-unveiling-key}, demonstrates that allocating intermediate computation enhances reasoning capacity even without meaningful semantics~\citep{pfau2024let, shi2025meaninglesstokensmeaningfulgains}.
Building on this, the field has shifted towards Latent Chain-of-Thought, which fully internalizes reasoning into hidden states~\citep{hao2025training, shen-etal-2025-codi}, with variants exploring parallelization~\citep{wu2025parallelcontinuouschainofthoughtjacobi}, test-time scaling~\citep{pttslatent} and kv-cache distillation~\citep{kava}. Complementing these horizontal structures, recursive and looped architectures have also emerged to scale reasoning depth through iterative state~\citep{geiping2025scaling, wang2025hierarchicalreasoningmodel, jolicoeurmartineau2025morerecursivereasoningtiny, zhu2025scalinglatentreasoninglooped}.
% However, supervision paradigms for these opaque states have bifurcated: Geometric Compression methods force latent states to approximate average embeddings~\citep{ccot, colar}, whereas Semantic Reconstruction approaches employ auxiliary decoding objectives to ensure the recoverability of logical steps~\citep{simcot, he2025semcot}. 
A central challenge in Latent CoT is how to supervise opaque continuous states.
Existing approaches can be broadly viewed as forms of space supervision.
Geometric compression (GC) constrains latent states to match target representations, such as average embeddings of reasoning steps~\citep{ccot, colar}.
Generative reconstruction (GR), in contrast, uses an auxiliary decoding objective to preserve the recoverability of explicit reasoning steps from latent states~\citep{simcot, he2025semcot}.
This distinction connects Latent CoT to broader representation learning objectives.
GC is analogous to JEPA-style representation prediction, where models are trained to predict target representations in latent space rather than reconstruct raw observations~\citep{lecun2022path}.
GR is closer in spirit to reconstruction-based objectives such as masked autoencoders, which recover the original input from a learned representation~\citep{he2022masked}.




\paragraph{Internal Dynamics and Theoretical Perspectives.}
~\citet{ton2025understanding} analyses characterize Explicit CoT as a potentially inefficient channel, where generated rationales often suffer from significant redundancy and low information gain.
Transcending the information bottleneck of discrete symbols, continuous representations offer provable advantages in efficiency: \citet{zou2025latentcollaborationmultiagentsystems} demonstrate that they achieve exponentially higher information transmission rates than discrete tokens. To address the resulting ``black box'' opacity, theoretical studies postulate that high-dimensional latent spaces enable \textit{reasoning by superposition}~\citep{zhu2025reasoning, xu2025formalcomparisonchainofthoughtlatent}. This mechanism allows models to process parallel logic paths simultaneously, a capability fundamentally unattainable by sequential discrete decoding~\citep{zhang2025soft}.
However, empirical scrutiny challenges this theoretical optimism: recent causal and adversarial analyses reveal that without robust grounding, latent tokens frequently degenerate into uninterpretable shortcuts or heuristic pattern-matching rather than encoding faithful step-by-step logic~\citep{lin-etal-2025-implicit, zhang2025latenttokensthinkcausal, yu2025llmsreallythinkstepbystep}, highlighting the critical difficulty in verifying the causal validity of internalized thought.
